\ifdefined\gkver\else\def\gkver{student}\fi
\documentclass[\gkver]{gaokaozhenti}
\newcommand{\ccnum}[1]{{\fontspec{Noto Sans CJK JP}#1}}
\begin{document}

\chapter{2014年四川理综}
%% paperType: 理综物理部分

\begin{enumerate}
\item
%% number: 1
%% paperName: 2014年四川理综第 1 题 6 分
%% typeId: 1
%% type: 单选题
%% chapter: 交流电
%% point: 变压器
%% method: 无
%% score: 6
%% degree: 650
%% duplicateId: 0
%% body:
如图所示，甲是远距离输电线路的示意图，乙是发电机输出电压随时间变化的图像，则\xzanswer{D}
\twopicture[labelA=2014四川理综01a, labelB=2014四川理综01b, numA=甲, numB=乙, labelprefix={}, wA=7.2cm, wB=5.8cm]
{figs/01a.png}
{figs/01b.png}
\fourchoices[answer=D]
{用户用电器上交流电的频率是 $100\UHz$}
{发电机输出交流电的电压有效值是 $500\UV$}
{输电线的电流只由降压变压器原副线圈的匝数比决定}
{当用户用电器的总电阻增大时，输电线上损失的功率减小}
%% answer: D
%% memo:
\memoanswer{
从图乙得到交流电的频率是 $50\UHz$，变压器在输电过程中不改变交流电的频率，A 错误；从图乙得到发电机输出电压的最大值是 $500\UV$，所以有效值为 $250\sqrt{2}\UV$，B 错误；输电线的电流是由降压变压器的负载电阻和输出电压决定的，C 错误；由于变压器的输出电压不变，当用户用电器的总电阻增大时，输出电流减小，根据电流与匝数成反比的关系可知，输电线上的电流减小，由 $P_{\text{线}}=I_{\text{线}}^{2}R_{\text{线}}$ 可知，输电线上损失的功率减小，选项 D 正确。
}

\item
%% number: 2
%% paperName: 2014年四川理综第 2 题 6 分
%% typeId: 1
%% type: 单选题
%% chapter: 光学
%% point: 光的电磁说
%% method: 无
%% score: 6
%% degree: 650
%% duplicateId: 0
%% body:
电磁波已广泛运用于很多领域。下列关于电磁波的说法符合实际的是\xzanswer{C}
\fourchoices[answer=C]
{电磁波不能产生衍射现象}
{常用的遥控器通过发出紫外线脉冲信号来遥控电视机}
{根据多普勒效应可以判断遥远天体相对于地球的运动速度}
{光在真空中运动的速度在不同惯性系中测得的数值可能不同}
%% answer: C
%% memo:
\memoanswer{
衍射现象是波的特有现象，A 错误；常用的遥控器通过发出红外线脉冲信号来遥控电视机，B 错误；遥远天体和地球的距离发生变化时，遥远天体的电磁波由于相对距离发生变化而出现多普勒效应，所以能测出遥远天体相对地球的运动速度，C 正确；光在真空中运动的速度在不同惯性系中测得的数值是相同的，即光速不变原理，D 错误。
}

\item
%% number: 3
%% paperName: 2014年四川理综第 3 题 6 分
%% typeId: 1
%% type: 单选题
%% chapter: 光学
%% point: 光的折射
%% method: 无
%% score: 6
%% degree: 650
%% duplicateId: 0
%% body:
如图所示，口径较大、充满水的薄壁圆柱形浅玻璃缸底有一发光小球，则\xzanswer{D}
\onepicture{\includesvg{figs/03}}
\fourchoices[answer=D]
{小球必须位于缸底中心才能从侧面看到小球}
{小球所发的光能从水面任何区域射出}
{小球所发的光从水中进入空气后频率变大}
{小球所发的光从水中进入空气后传播速度变大}
%% answer: D
%% memo:
\memoanswer{
光从水中进入空气，只要在没有发生全反射的区域，就可以看到光线射出，所以 A、B 错误；光的频率是由光源决定的，与介质无关，所以 C 错误；由 $v=\frac{c}{n}$ 得，光从水中进入空气后传播速度变大，所以 D 正确。
}

\item
%% number: 4
%% paperName: 2014年四川理综第 4 题 6 分
%% typeId: 1
%% type: 单选题
%% chapter: 直线运动
%% point: 运动的合成
%% method: 无
%% score: 6
%% degree: 650
%% duplicateId: 0
%% body:
有一条两岸平直、河水均匀流动、流速恒为 $v$ 的大河。小明驾着小船渡河，去程时船头指向始终与河岸垂直，回程时行驶路线与河岸垂直。去程与回程所用时间的比值为 $k$，船在静水中的速度大小相同，则小船在静水中的速度大小为\xzanswer{B}
\fourchoices[answer=B]
{$\frac{kv}{\sqrt{k^{2}-1}}$}
{$\frac{v}{\sqrt{1-k^{2}}}$}
{$\frac{kv}{\sqrt{1-k^{2}}}$}
{$\frac{v}{\sqrt{k^{2}-1}}$}
%% answer: B
%% memo:
\memoanswer{
设河岸宽为 $d$，船速为 $u$，则根据渡河时间关系得 $\frac{d}{u}:\frac{d}{\sqrt{u^{2}-v^{2}}}=k$，解得 $u=\frac{v}{\sqrt{1-k^{2}}}$，所以 B 选项正确。
}

\item
%% number: 5
%% paperName: 2014年四川理综第 5 题 6 分
%% typeId: 1
%% type: 单选题
%% chapter: 机械振动
%% point: 振动图象
%% method: 无
%% score: 6
%% degree: 450
%% duplicateId: 0
%% body:
如图所示，甲为 $t=1\Us$ 时某横波的波形图像，乙为该波传播方向上某一质点的振动图像，距该质点 $\Delta x=0.5\Um$ 处质点的振动图像可能是\xzanswer{A}
\twopicture[labelA=2014四川理综05a, labelB=2014四川理综05b, numA=甲, numB=乙, labelprefix={}, wA=5.4cm, wB=4.6cm]
{figs/05a.png}
{figs/05b.png}
\fourchoices[answer=A, ispicture=true, h=2.0cm]
{figs/05c.png}
{figs/05d.png}
{figs/05e.png}
{figs/05f.png}
%% answer: A
%% memo:
\memoanswer{
从甲图可以得到波长为 $2\Um$，从乙图可以得到周期为 $2\Us$，即波速为 $1\Ums$；由乙图的振动图像可以找到 $t=1\Us$ 时，该质点位移为负，并且向下运动，距该质点 $\Delta x=0.5\Um$ 处的质点与该质点的振动情况相差 $\frac{T}{4}$，即将乙图中的图像向左或右平移 $\frac{1}{4}$ 周期即可得到距该质点 $\Delta x=0.5\Um$ 处质点的振动图像，故只有 A 正确。
}

\item
%% number: 6
%% paperName: 2014年四川理综第 6 题 6 分
%% typeId: 2
%% type: 多选题
%% chapter: 电磁感应
%% point: 电磁感应受力分析
%% method: 无
%% score: 6
%% degree: 450
%% duplicateId: 0
%% body:
如图所示，一不计电阻的光滑 U 形金属框水平放置，光滑、竖直玻璃挡板 H、P 固定在框上，H、P 的间距很小。质量为 $0.2\Ukg$ 的细金属杆 CD 恰好无挤压地放在两挡板之间，与金属框接触良好并围成边长为 $1\Um$ 的正方形，其有效电阻为 $0.1\UO$。此时在整个空间加方向与水平面成 $30^\circ$ 角且与金属杆垂直的匀强磁场，磁感应强度随时间变化规律是 $B=(0.4-0.2t)\UT$，图示磁场方向为正方向。框、挡板和杆不计形变。则\xzanswer{AC}
\onepicture{\includesvg{figs/06}}
\fourchoices[answer=AC]
{$t=1\Us$ 时，金属杆中感应电流方向从 C 到 D}
{$t=3\Us$ 时，金属杆中感应电流方向从 D 到 C}
{$t=1\Us$ 时，金属杆对挡板 P 的压力大小为 $0.1\UN$}
{$t=3\Us$ 时，金属杆对挡板 H 的压力大小为 $0.2\UN$}
%% answer: AC
%% memo:
\memoanswer{
由于 $B=(0.4-0.2t)\UT$，在 $t=1\Us$ 时穿过平面的磁通量向下并减少，则根据楞次定律可以判断，金属杆中感应电流方向从 C 到 D，A 正确。在 $t=3\Us$ 时穿过平面的磁通量向上并增加，则根据楞次定律可以判断，金属杆中感应电流方向仍然是从 C 到 D，B 错误。由法拉第电磁感应定律得 $E=\frac{\Delta\Phi}{\Delta t}=\frac{\Delta B}{\Delta t}S\sin 30^\circ=0.1\UV$，由闭合电路的欧姆定律得电路电流 $I=\frac{E}{R}=1\UA$，在 $t=1\Us$ 时，$B=0.2\UT$，方向斜向下，电流方向从 C 到 D，金属杆对挡板 P 的压力水平向右，大小为 $F_{P}=BIL\sin 30^\circ=0.1\UN$，C 正确。同理，在 $t=3\Us$ 时，金属杆对挡板 H 的压力水平向左，大小为 $F_{H}=BIL\sin 30^\circ=0.1\UN$，D 错误。
}

\item
%% number: 7
%% paperName: 2014年四川理综第 7 题 6 分
%% typeId: 2
%% type: 多选题
%% chapter: 运动和力
%% point: 牛顿定律的应用
%% method: 无
%% score: 6
%% degree: 450
%% duplicateId: 0
%% body:
如右图所示，水平传送带以速度 $v_{1}$ 匀速运动，小物体 P、Q 由通过定滑轮且不可伸长的轻绳相连，$t=0$ 时刻 P 在传送带左端具有速度 $v_{2}$，P 与定滑轮间的绳水平，$t=t_{0}$ 时刻 P 离开传送带。不计定滑轮质量和摩擦，绳足够长。正确描述小物体 P 速度随时间变化的图像可能是\xzanswer{BC}
\onepicture[align=right]{\includesvg{figs/07a}}
\fourchoices[answer=BC, ispicture=true, h=2.2cm]
{figs/07b.png}
{figs/07c.png}
{figs/07d.png}
{figs/07e.png}
%% answer: BC
%% memo:
\memoanswer{
若 P 在传送带左端时的速度 $v_{2}$ 小于 $v_{1}$，则 P 受到向右的摩擦力，当 P 受到的摩擦力大于绳的拉力时，P 做加速运动，则有两种可能：第一种是一直做加速运动，第二种是先做加速运动，当速度达到 $v_{1}$ 后做匀速运动，所以 B 正确；当 P 受到的摩擦力小于绳的拉力时，P 做减速运动，也有两种可能：第一种是一直做减速运动，从右端滑出；第二种是先做减速运动再做反向加速运动，从左端滑出。若 P 在传送带左端具有的速度 $v_{2}$ 大于 $v_{1}$，则小物体 P 受到向左的摩擦力，使 P 做减速运动，则有三种可能：第一种是一直做减速运动，第二种是速度先减到 $v_{1}$，之后若 P 受到绳的拉力和静摩擦力作用而处于平衡状态，则其以速度 $v_{1}$ 做匀速运动，第三种是速度先减到 $v_{1}$，之后若 P 所受的静摩擦力小于绳的拉力，则 P 将继续减速直到速度减为 0，再反向做加速运动并且摩擦力反向，加速度不变，从左端滑出，所以 C 正确。
}

\item
%% number: 8
%% paperName: 2014年四川理综第 8 题 11 分
%% typeId: 4
%% type: 实验题
%% chapter: 电路
%% point: 电路实验
%% method: 无
%% score: 11
%% degree: 350
%% duplicateId: 0
%% body:
（1）小文同学在探究物体做曲线运动的条件时，将一条形磁铁放在桌面的不同位置，让小钢珠在水平桌面上从同一位置以相同初速 $v_{0}$ 运动，得到不同轨迹。图中 a、b、c、d 为其中四条运动轨迹，磁铁放在位置 A 时，小钢珠的运动轨迹是\_\_\_\_\_\_\_\_\_（填轨迹字母代号），磁铁放在位置 B 时，小钢珠的运动轨迹是\_\_\_\_\_\_\_\_\_（填轨迹字母代号）。实验表明，当物体所受合外力的方向跟它的速度方向\_\_\_\_\_\_\_\_\_\_（选填“在”或“不在”）同一直线上时，物体做曲线运动。

\onepicture[align=right]{\includesvg{figs/08}}

（2）右图是测量阻值约几十欧的未知电阻 $R_{x}$ 的原理图，图中 $R_{0}$ 是保护电阻（$10\UO$），$R_{1}$ 是电阻箱（$0\sim 99.9\UO$），$R$ 是滑动变阻器，A$_{1}$ 和 A$_{2}$ 是电流表，$E$ 是电源（电动势 $10\UV$，内阻很小）。

在保证安全和满足要求的情况下，使测量范围尽可能大。实验具体步骤如下：

（i）连接好电路，将滑动变阻器 $R$ 调到最大；

（ii）闭合 S，从最大值开始调节电阻箱 $R_{1}$，先调 $R_{1}$ 为适当位置，再调节滑动变阻器 $R$，使 A$_{1}$ 示数 $I_{1}=0.15\UA$，记下此时电阻箱的阻值 $R_{1}$ 和 A$_{2}$ 的示数 $I_{2}$；

（iii）重复步骤（ii），再测量 6 组 $R_{1}$ 和 $I_{2}$ 值；

（iv）将实验测得的 7 组数据在坐标纸上描点。

根据实验回答以下问题：

\begin{enumerate}
	\item
	现有四只供选用的电流表：
	（A）电流表（$0\sim 3\UmA$，内阻为 $2.0\UO$）
	（B）电流表（$0\sim 3\UmA$，内阻未知）
	（C）电流表（$0\sim 0.3\UA$，内阻为 $5.0\UO$）
	（D）电流表（$0\sim 0.3\UA$，内阻未知）
	A$_{1}$ 应选用\tkanswer[1.5cm]{D}，A$_{2}$ 应选\tkanswer[1.5cm]{C}。

	\item
	测得一组 $R_{1}$ 和 $I_{2}$ 值后，调整电阻箱 $R_{1}$，使其阻值变小，要使 A$_{1}$ 示数 $I_{1}=0.15\UA$，应让滑动变阻器 $R$ 接入电路的阻值\tkanswer[1.8cm]{变大}（选填“不变”、“变大”或“变小”）。

	\item
	在坐标纸上画出 $R_{1}$ 与 $I_{2}$ 的关系图。
	\onepicture[w=5.5cm]{\includesvg{figs/08b}}

	\item
	根据以上实验得出 $R_{x}=$\tkanswer[1.5cm]{31} $\UO$。
\end{enumerate}
\onepicture[align=right]{\includesvg{figs/08a}}
%% answer: 
\jdanswer{
\textbf{（1）}b、c，不在

\textbf{（2）}
\begin{enumerate}
	\item
	D、C

	\item
	变大

	\item
	关系图线如图
	\onepicture[w=5.5cm]{\includesvg{figs/08_an}}

	\item
	$31\UO$

\end{enumerate}
}
%% memo:
\memoanswer{
（1）由磁铁的引力作用及曲线运动的规律可知 A 处时沿 b 轨迹，B 处时沿 c 轨迹。

（2）\ccnum{①}A$_{1}$ 的示数能达到 $0.15\UA$，A$_{2}$ 的示数由图像可知能达到 $0.3\UA$，故 A$_{1}$、A$_{2}$ 的量程均选 $0.3\UA$，由电路图可列出关系式 $(R_{x}+R_{A2})I_{2}=(R_{0}+R_{1}+R_{A1})I_{1}$，整理后可得 $R_{A2}+R_{x}=\frac{I_{1}}{I_{2}}(R_{0}+R_{1}+R_{A1})$，由此可知，若 $R_{A1}$ 已知，则无论 $R_{1}$、$I_{2}$ 如何变化，$R_{x}+R_{A2}$ 均为定值，无法得到 $R_{x}$，故应使 $R_{A2}$ 已知，即 A$_{1}$ 选 D，A$_{2}$ 选 C。

\ccnum{②}当 $R_{1}$ 减小时，如果在滑动变阻器的电阻值保持不变的情况下，电路的总电阻减小，由闭合电路的欧姆定律可得总电流 $I$ 增大，由分压关系知，并联部分分得的电压减小，则 $I_{2}$ 减小，由 $I_{1}=I-I_{2}$ 得 $I_{1}$ 增大，要使 $I_{1}=0.15\UA$，则需滑动变阻器分得的电压增大，即 $R$ 的阻值变大。

\ccnum{③}在坐标纸上描出各数据点，作出 $R_{1}-I_{2}$ 关系图线，如图（答）所示。

\ccnum{④}根据 $(R_{x}+R_{A2})I_{2}=(R_{0}+R_{1}+R_{A1})I_{1}$，可得 $R_{1}=\frac{R_{x}+R_{A2}}{I_{1}}I_{2}-(R_{0}-R_{A1})$，即 $R_{1}-I_{2}$ 图像的斜率 $k=\frac{R_{x}+R_{A2}}{I_{1}}$，根据图像并代入相关数据，可得 $R_{x}=31\UO$。
}

\item
%% number: 9
%% paperName: 2014年四川理综第 9 题 15 分
%% typeId: 6
%% type: 计算题
%% chapter: 曲线运动
%% point: 人造地球卫星
%% method: 无
%% score: 15
%% degree: 650
%% duplicateId: 0
%% body:
石墨烯是近些年发现的一种新材料，其超高强度及超强导电、导热等非凡的物理化学性质有望使 21 世纪的世界发生革命性的变化，其发现者由此获得 2010 年诺贝尔物理学奖。用石墨烯制作超级缆绳，人类搭建“太空电梯”的梦想有望在本世纪实现。科学家们设想，通过地球同步轨道站向地面垂下一条缆绳至赤道基站，电梯仓沿着这条缆绳运行，实现外太空和地球之间便捷的物资交换。
\begin{enumerate}
	\item
	若“太空电梯”将货物从赤道基站运到距离地面高度为 $h_{1}$ 的同步轨道站，求轨道站内质量为 $m_{1}$ 的货物相对地心运动的动能。设地球自转角速度为 $\omega$，地球半径为 $R$。
	\item
	当电梯仓停在距地面高度 $h_{2}=4R$ 的站点时，求仓内质量 $m_{2}=50\Ukg$ 的人对水平地板的压力大小。取地面附近重力加速度 $g=10\Umsq$，地球自转角速度 $\omega=7.3\times10^{-5}\Urads$，地球半径 $R=6.4\times10^{3}\Ukm$。
\end{enumerate}
\onepicture[align=right]{figs/09.jpg}
%% answer: 
\jdanswer{
\begin{enumerate}
	\item
	$\frac{1}{2}m_{1}\omega^{2}(R+h_{1})^{2}$

	\item
	$11.5\UN$

\end{enumerate}
}
%% memo:
\memoanswer{
（1）设货物相对地心的距离为 $r_{1}$，线速度为 $v_{1}$，则

$r_{1}=R+h_{1}$\ccnum{①}

$v_{1}=r_{1}\omega$\ccnum{②}

货物对地心的动能为 $E_{k}=\frac{1}{2}m_{1}v_{1}^{2}$\ccnum{③}

联立\ccnum{①}\ccnum{②}\ccnum{③}式得 $E_{k}=\frac{1}{2}m_{1}\omega^{2}(R+h_{1})^{2}$\ccnum{④}

说明：\ccnum{①}\ccnum{②}\ccnum{③}\ccnum{④}式各 1 分。

（2）设地球质量为 $M$，人相对地心的距离为 $r_{2}$，向心加速度为 $a_{\text{向}}$，受地球的万有引力为 $F$，则

$r_{2}=R+h_{2}$\ccnum{⑤}

$a_{\text{向}}=\omega^{2}r_{2}$\ccnum{⑥}

$F=G\frac{m_{2}M}{r_{2}^{2}}$\ccnum{⑦}

$g=\frac{GM}{R^{2}}$\ccnum{⑧}

设水平地板对人的支持力大小为 $N$，人对水平地板的压力大小为 $N'$，则

$F-N=m_{2}a_{\text{向}}$\ccnum{⑨}

$N'=N$\ccnum{⑩}

联立\ccnum{⑤}～\ccnum{⑩}式并代入数据得 $N'=11.5\UN$\ccnum{⑪}

说明：\ccnum{⑥}\ccnum{⑦}\ccnum{⑧}\ccnum{⑨}式各 2 分，\ccnum{⑤}\ccnum{⑩}\ccnum{⑪}式各 1 分。
}

\item
%% number: 10
%% paperName: 2014年四川理综第 10 题 17 分
%% typeId: 6
%% type: 计算题
%% chapter: 磁场
%% point: 带电粒子在复合场中的运动
%% method: 无
%% score: 17
%% degree: 450
%% duplicateId: 0
%% body:
在如图所示的竖直平面内，水平轨道 CD 和倾斜轨道 GH 与半径 $r=\frac{9}{44}\Um$ 的光滑圆弧轨道分别相切于 D 点和 G 点，GH 与水平面的夹角 $\theta=37^\circ$。过 G 点、垂直于纸面的竖直平面左侧有匀强磁场，磁场方向垂直于纸面向里，磁感应强度 $B=1.25\UT$；过 D 点、垂直于纸面的竖直平面右侧有匀强电场，电场方向水平向右，电场强度 $E=1\times10^{4}\UNC$。小物体 $P_{1}$ 质量 $m=2\times10^{-3}\Ukg$、电荷量 $q=+8\times10^{-6}\UC$，受到水平向右的推力 $F=9.98\times10^{-3}\UN$ 的作用，沿 CD 向右做匀速直线运动，到达 D 点后撤去推力。当 $P_{1}$ 到达倾斜轨道底端 G 点时，不带电的小物体 $P_{2}$ 在 GH 顶端静止释放，经过时间 $t=0.1\Us$ 与 $P_{1}$ 相遇。$P_{1}$ 与 $P_{2}$ 与轨道 CD、GH 间的动摩擦因数均为 $\mu=0.5$，取 $g=10\Umsq$，$\sin 37^\circ=0.6$，$\cos 37^\circ=0.8$，物体电荷量保持不变，不计空气阻力。求：
\begin{enumerate}
	\item
	小物体 $P_{1}$ 在水平轨道 CD 上运动速度 $v$ 的大小；
	\item
	倾斜轨道 GH 的长度 $s$。
\end{enumerate}
\onepicture[align=right]{\includesvg{figs/10}}
%% answer: 
\jdanswer{
\begin{enumerate}
	\item
	$4\Ums$

	\item
	$0.56\Um$

\end{enumerate}
}
%% memo:
\memoanswer{
（1）设小物体 $P_{1}$ 在匀强磁场中运动的速度为 $v$，受到向上的洛伦兹力为 $F_{1}$，受到的摩擦力为 $f$，则

$F_{1}=qvB$\ccnum{①}

$f=\mu(mg-F_{1})$\ccnum{②}

由题意，水平方向合力为零

$F-f=0$\ccnum{③}

联立\ccnum{①}\ccnum{②}\ccnum{③}式，代入数据解得

$v=4\Ums$\ccnum{④}

说明：\ccnum{①}\ccnum{③}式各 1 分，\ccnum{②}\ccnum{④}式各 2 分。

（2）设 $P_{1}$ 在 G 点的速度大小为 $v_{G}$，由于洛伦兹力不做功，根据动能定理

$qEr\sin\theta-mgr(1-\cos\theta)=\frac{1}{2}mv_{G}^{2}-\frac{1}{2}mv^{2}$\ccnum{⑤}

$P_{1}$ 在 GH 上运动，受到重力、电场力和摩擦力的作用，设加速度为 $a_{1}$，根据牛顿第二定律

$qE\cos\theta-mg\sin\theta-\mu(mg\cos\theta+qE\sin\theta)=ma_{1}$\ccnum{⑥}

$P_{1}$ 与 $P_{2}$ 在 GH 上相遇时，设 $P_{1}$ 在 GH 运动的距离为 $s_{1}$，则

$s_{1}=v_{G}t+\frac{1}{2}a_{1}t^{2}$\ccnum{⑦}

设 $P_{2}$ 质量为 $m_{2}$，在 GH 上运动的加速度为 $a_{2}$，则

$m_{2}g\sin\theta-\mu m_{2}g\cos\theta=m_{2}a_{2}$\ccnum{⑧}

$P_{1}$ 与 $P_{2}$ 在 GH 上相遇时，设 $P_{2}$ 在 GH 运动的距离为 $s_{2}$，则

$s_{2}=\frac{1}{2}a_{2}t^{2}$\ccnum{⑨}

联立\ccnum{⑤}～\ccnum{⑨}式，代入数据得

$s=s_{1}+s_{2}$\ccnum{⑩}

$s=0.56\Um$\ccnum{⑪}

说明：\ccnum{⑦}\ccnum{⑧}\ccnum{⑨}\ccnum{⑩}式各 1 分，\ccnum{⑥}\ccnum{⑪}式各 2 分。
}

\item
%% number: 11
%% paperName: 2014年四川理综第 11 题 19 分
%% typeId: 6
%% type: 计算题
%% chapter: 磁场
%% point: 带电粒子在复合场中的运动
%% method: 无
%% score: 19
%% degree: 250
%% duplicateId: 0
%% body:
如图所示，水平放置的不带电的平行金属板 p 和 b 相距 $h$，与图示电路相连，金属板厚度不计，忽略边缘效应。p 板上表面光滑，涂有绝缘层，其上 O 点右侧相距 $h$ 处有小孔 K；b 板上有小孔 T，且 O、T 在同一条竖直线上，图示平面为竖直平面。质量为 $m$、电荷量为 $-q$（$q>0$）的静止粒子被发射装置（图中未画出）从 O 点发射，沿 p 板上表面运动时间 $t$ 后到达 K 孔，不与板碰撞地进入两板之间。粒子视为质点，在图示平面内运动，电荷量保持不变，不计空气阻力，重力加速度大小为 $g$。
\begin{enumerate}
	\item
	求发射装置对粒子做的功；
	\item
	电路中的直流电源内阻为 $r$，开关 S 接“1”位置时，进入板间的粒子落在 b 板上的 A 点，A 点与过 K 孔竖直线的距离为 $l$。此后将开关 S 接“2”位置，求阻值为 $R$ 的电阻中的电流强度；
	\item
	若选用恰当直流电源，电路中开关 S 接“1”位置，使进入板间的粒子受力平衡，此时在板间某区域加上方向垂直于图面的、磁感应强度大小合适的匀强磁场（磁感应强度 $B$ 只能在 $0\sim B_{m}=\frac{(\sqrt{21}+5)m}{(\sqrt{21}-2)qt}$ 范围内选取），使粒子恰好从 b 板的 T 孔飞出，求粒子飞出时速度方向与 b 板面的夹角的所有可能值（可用反三角函数表示）。
\end{enumerate}
\onepicture[align=right]{\includesvg{figs/11}}
%% answer: 
\jdanswer{
\begin{enumerate}
	\item
	$\frac{mh^{2}}{2t^{2}}$

	\item
	$\frac{mh}{q(R+r)}\left(g-\frac{2h^{3}}{l^{2}t^{2}}\right)$

	\item
	$0<\theta\leqslant\arcsin\frac{2}{5}$

\end{enumerate}
}
%% memo:
\memoanswer{
（1）设粒子在 p 板上做匀速直线运动的速度为 $v_{0}$，有

$h=v_{0}t$\ccnum{①}

设发射装置对粒子做的功为 $W$，由动能定理

$W=\frac{1}{2}mv_{0}^{2}$\ccnum{②}

联立\ccnum{①}\ccnum{②}式可得 $W=\frac{mh^{2}}{2t^{2}}$\ccnum{③}

说明：\ccnum{①}\ccnum{②}式各 2 分，\ccnum{③}式 1 分。

（2）S 接“1”位置时，电源的电动势 $E_{0}$ 与板间电势差 $U$ 有

$E_{0}=U$\ccnum{④}

板间产生匀强电场的场强为 $E$，粒子进入板间时有水平方向的速度 $v_{0}$，在板间受到竖直方向的重力和电场力作用而做类平抛运动，设加速度为 $a$，运动时间为 $t_{1}$，有

$U=Eh$\ccnum{⑤}

$mg-qE=ma$\ccnum{⑥}

$h=\frac{1}{2}at_{1}^{2}$\ccnum{⑦}

$l=v_{0}t_{1}$\ccnum{⑧}

S 接“2”位置时，则在电阻 $R$ 上流过的电流 $I$ 满足

$I=\frac{E_{0}}{R+r}$\ccnum{⑨}

联立\ccnum{①}\ccnum{④}～\ccnum{⑨}式得

$I=\frac{mh}{q(R+r)}\left(g-\frac{2h^{3}}{l^{2}t^{2}}\right)$\ccnum{⑩}

说明：\ccnum{④}～\ccnum{⑩}式各 1 分。

（3）由题意知此时在板间运动的粒子重力与电场力平衡，当粒子从 K 进入板间后立即进入磁场做匀速圆周运动，如图所示，粒子从 D 点出磁场区域后沿 DT 做匀速直线运动，DT 与 b 板上表面的夹角为题目所求夹角 $\theta$，磁场的磁感应强度 $B$ 取最大值时的夹角 $\theta$ 为最大值 $\theta_{m}$，设粒子做匀速圆周运动的半径为 $R$，有

\onepicture[align=right, w=7.5cm]{figs/11_an.png}

$qv_{0}B=\frac{mv_{0}^{2}}{R}$\ccnum{⑪}

过 D 点作 b 板的垂线与 b 板的上表面交于 G，由几何关系有

$\overline{DG}=h-R(1+\cos\theta)$\ccnum{⑫}

$\overline{TG}=h+R\sin\theta$\ccnum{⑬}

$\tan\theta=\frac{\sin\theta}{\cos\theta}=\frac{\overline{DG}}{\overline{TG}}$\ccnum{⑭}

联立\ccnum{①}\ccnum{⑪}～\ccnum{⑭}式，将 $B=B_{m}$ 代入，求得

$\theta_{m}=\arcsin\frac{2}{5}$\ccnum{⑮}

当 $B$ 逐渐减小，粒子做匀速圆周运动的半径为 $R$ 也随之变大，D 点向 b 板靠近，DT 与 b 板上表面的夹角 $\theta$ 也越变越小，当 D 点无限接近于 b 板上表面时，粒子离开磁场后在板间几乎沿着 b 板上表面从 T 孔飞出板间区域，此时 $B_{m}>B>0$ 满足题目要求，夹角 $\theta$ 趋近 $\theta_{0}$，即

$\theta_{0}=0$\ccnum{⑯}

则题目所求为

$0<\theta\leqslant\arcsin\frac{2}{5}$\ccnum{⑰}

说明：\ccnum{⑪}～\ccnum{⑰}式各 1 分。
}

\end{enumerate}

\end{document}
